\documentclass[11pt]{article}
\usepackage[margin=1in]{geometry}
\usepackage{titlesec}
\usepackage{nopageno}

% Section formatting for NSF Project Summary
\titleformat{\section}{\normalfont\bfseries}{\thesection}{1em}{}
\titlespacing*{\section}{0pt}{0.5\baselineskip}{0.5\baselineskip}

\begin{document}

\noindent \textbf{Project Summary: Phase I CAMEL-CN: Modeling the Messy}

\section*{Overview}
The Phase I CAMEL Collaborative Network (CAMEL-CN) will address the lack of K-12 educational datasets documenting how learners navigate real-world data complexity, a foundational capability for AI literacy. This Novel Dataset Pathway (1a) project immerses 4,500 high school students (Grades 9--12) in authentic industrial data contexts, specifically AI infrastructure cooling systems and advanced materials processing, to generate a multi-layered, AI-ready Student Math Modeling Learning Dataset. By integrating researchers in the science of learning, education practitioners from the Center for Science and the Schools (CSATS), and disciplinary data scientists, this network bridges fundamental knowledge of human cognition with classroom practice. The project will use a unique ``Data Translation'' framework to convert high-dimensional engineering research data into pedagogical materials with adjustable ``complexity dials.'' This allows students to encounter the statistical, structural, and provenance challenges typical of STEM careers while remaining appropriately scaffolded for high school Algebra and Statistics standards.

\section*{Intellectual Merit}
No existing educational dataset captures how students reason through authentic, messy industrial data in open-ended computational modeling tasks. Resources such as ASSISTments, PSLC DataShop, TIMSS, and NAEP document whether students reached correct answers, but treat the reasoning process as a black box, none capture the deleted attempts, debugging cycles, and iterative refinements that reveal how mathematical modeling competencies actually develop. This project fills that gap directly.
The intellectual merit lies in the convergence of learning sciences, data science, and classroom practice to investigate how students develop mathematical modeling competencies when forced to reconcile sensor noise, high dimensionality, and data provenance gaps,  the challenges traditionally sanitized from curricula. The network will establish a replicable methodological infrastructure: a Data Translation Framework for converting high-dimensional engineering research data into pedagogically scaffolded materials, and AI-ready annotation protocols aligned with xAPI and Caliper standards. Process-level telemetry captured via Jupyter/IPython kernel hooks will map iterative student reasoning trajectories at scale, enabling study of the behavioral signatures of productive struggle and identification of the cognitive thresholds where data-driven friction becomes unproductive. The resulting FAIR-compliant dataset, tens of millions of annotated events from 4,500 students, will provide the scale and structure required to train next-generation intelligent tutoring systems and formative assessment tools optimized for authentic data contexts.

\section*{Broader Impacts}
The broader impacts are designed to create a recursive feedback loop for workforce development by preparing students for data-intensive STEM careers. The project will explicitly address current sampling deficiencies by prioritizing rural Pennsylvania school districts and Title I schools, populations frequently underrepresented in large-scale learning analytics. Through a Research-Practice Partnership (RPP) model, rural math teachers will be empowered as co-researchers, participating in data annotation and analysis during content-intensive Summer Institutes. This will build long-term STEM capacity in isolated districts. To promote usage, the network will publish the dataset on established repositories like ICPSR and the Open Science Framework, supplemented by tutorial notebooks and analysis scripts to lower the barrier for other researchers to adopt process-level data. Finally, coordination with the Phase II National Coordinator will ensure these outcomes are tracked and integrated into a national infrastructure for advancing mathematics education.

\end{document}